\begin{lemma}[\textup{LEM\_Stable}: recoloring chains outside a good hierarchy]
Let $T$ be a positive-rung ETT satisfying maximum property and carrying a
good hierarchy under its generating coloring $\varphi_n$.  Let
$\sigma_n$ be obtained by recoloring finitely many
$(\alpha,\beta)$-chains fully contained in $G-V(T)$.  Keep $G,k,e_0$
fixed but allow the certified input coloring to vary.  Assume that for
every positive-rung ETT $T'$ with arbitrary certified input coloring,
maximum property, and rung number $s<n$, the following six assertions
hold.
\begin{enumerate}
\item $V(T')$ is elementary under its generating coloring
$\varphi'_s$.
\item Under every proper $(T'_{s+1},D'_s,\varphi'_s)$-stable coloring
leaving exactly $e_0$ uncolored, $T'_{s+1}$ is interchangeable for any
two colors of which at least one is missing on it.
\item At every PE slot, stability descends to the prefix before the
support and to the ladder stage indexed by the support's first
appearance; stored pairs at distinct PE slots with distinct supports
and equal first-appearance indices are disjoint.
\item At a terminal PE, the stored two-color component through its
support meets the terminal ladder tree only at that support.
\item Every nonminimum tree endpoint of a defective boundary color
precedes each later PE support or SE extension vertex permitted by its
first-appearance index.
\item Every proper terminal stable coloring leaving exactly $e_0$
uncolored is a $\varphi'_s\bmod T'_s$ coloring.
\end{enumerate}
Then:
\begin{enumerate}
\item $\sigma_n$ is proper, leaves exactly $e_0$ uncolored, is
$(T,D_n,\varphi_n)$-stable, and is $(T,\varphi_n)$-invariant;
\item the same terminal tree carries a corresponding ETT certificate
under $\sigma_n$ and satisfies maximum property;
\item the same ordered stages and dividers form a good hierarchy under
$\sigma_n$, with the same $D_{n,j}$ and the same $\Gamma^j_\eta$ for
every $\eta\in D_{n,j}$; and
\item if $T$ is
$(\bigcup_{\eta\in D_{n,q+1}}\Gamma^q_\eta)^{-}$-closed under
$\varphi_n$, it is so under $\sigma_n$.
\end{enumerate}
\end{lemma}
\begin{proof}
Write $n$ for the terminal rung, $T_n$ for the terminal ladder tree,
$T_{n+1}$ for the last tree of the certified series, and
$D_n$ for the terminal stored-color set.  Apply
\noderef{HierarchyOutsideKempeStable} with the given ETT, hierarchy base,
hierarchy, properness, one-uncolored-edge hypothesis, and outside
recoloring.  It gives that $\sigma_n$ is proper, leaves exactly $e_0$
uncolored, is $(T,D_n,\varphi_n)$-stable, and is
$(T,\varphi_n)$-invariant.  It also gives
$(T_n\oplus R_n,D_n,\varphi_n)$-stability whenever the terminal kind is
PE.  This proves every assertion in (i).

We shall need the stronger pointwise form of the outside-recoloring
observation, so we obtain it independently rather than reading it from
the preceding lemma's proof.  Apply
\noderef{OutsideKempeRecoloringStable} directly, with control set $D_n$,
to the stated proper coloring, the fact that exactly $e_0$ is uncolored,
and the given outside recoloring.  Its last conclusion gives
\[
 \sigma_n(f)=\varphi_n(f)
 \quad\text{for every edge identity $f$ incident with $T$.}   \tag{1}
\]
Its stable-coloring conclusion, expanded using the vertex clause of
\noderef{StableColoring}, also gives
\[
 \overline{\sigma_n}(v)=\overline{\varphi_n}(v)
 \quad\text{for every $v\in V(T)$.}                           \tag{2}
\]

We next produce the corresponding ETT certificate required in (ii).
By \noderef{ExtendedTashkinovTree}, $T_n$ is an ordered initial segment
of $T$, its length is strictly smaller than that of $T$, $T$ is an
ordered initial segment of $T_{n+1}$, and the generating coloring is the
terminal series coloring $\varphi_n$.  The entrywise definition in
\noderef{TreeSequenceInitialSegment}, together with
\noderef{TreeSequenceVertexSet}, therefore gives
$V(T_n)\subseteq V(T)$.  Equations (1) and (2) restrict to the edge and
vertex clauses of $(T_n,D_n,\varphi_n)$-stability: for the edge clause,
an edge incident with $T_n$ is incident with the same vertex in $T$ by
\noderef{TreeIncidentEdge}; for the vertex clause, use the displayed
vertex inclusion.  Apply \noderef{Td} to this restricted stability,
the given lower-rung package, positivity, and maximum property.  Its
conclusion is
\[
                 \sigma_n\text{ is a }\varphi_n\bmod T_n
                 \text{ coloring}.                           \tag{3}
\]

It remains to verify, rather than merely assert, that $T$ is eligible
for the universal clause of \noderef{ModColoring}.  The final certificate
in \noderef{TashkinovSeries} is a
\noderef{TashkinovSeriesStep}.  Positivity of $n$ excludes its initial
constructor.  In the PE constructor, its closure field is a
\noderef{TashkinovClosure} of $T_n$ under $\varphi_n$ whose last tree is
$T_{n+1}$.  In either the SE or RE constructor, it supplies an individual
edge identity $f_n$, the exact equality $F_n=\{f_n\}$, a tree $A$ with
\[
              A=T_n+f_n
              \quad\text{in the sense of }
              \noderef{TreeSequenceAdjoinEdge},               \tag{4}
\]
and a $\varphi_n$-closure from $A$ to $T_{n+1}$.

We first truncate the appropriate old closure at the displayed tree
$T$.  Expanding \noderef{TashkinovClosure}, its first field is a finite
\noderef{TashkinovAugmentationSequence}.  Each step in that sequence is
a \noderef{TashkinovAugmentation}; hence it increases the length by
exactly one and retains the preceding tree as the exact ordered initial
segment.  In the PE case the seed $T_n$ is strictly shorter than $T$.
In the SE/RE case (4) has length $|T_n|+1\leq |T|$, because $T$ strictly
extends $T_n$.  Moreover, both $A$ and $T$ are initial segments of
$T_{n+1}$: composing the one-step initial-segment fields in the SE/RE
closure sequence gives this for $A$, and the ETT certificate gives it for
$T$.  Their entries agree through the shorter length by
\noderef{TreeSequenceInitialSegment}, so $A$ is an initial segment of
$T$.

The augmentation lengths in the closure sequence consequently pass
through the length of $T$.  The member having that length and $T$ are
both ordered initial segments of $T_{n+1}$ of the same length; equality
of every vertex and edge entry makes them the same ordered tree-sequence.
Taking the initial part of the finite sequence therefore gives
\[
 \begin{cases}
  T_n\longrightarrow T,&\text{in the PE case},\\
  T_n+f_n=A\longrightarrow T,&\text{in the SE/RE cases},
 \end{cases}                                                   \tag{5}
\]
as augmentation sequences under $\varphi_n$.  The second route retains
the exact edge identity and singleton equality in (4), not just its
endpoints.

We now transport every step of (5) from $\varphi_n$ to $\sigma_n$.
Consider one step $A_i\to A_{i+1}$.  Its new last edge $g_i$ joins a
vertex of $A_i$ to the new vertex by
\noderef{TashkinovAugmentation}.  Since $A_i$ is an initial segment of
$T$, that old endpoint lies in $V(T)$ by
\noderef{TreeSequenceVertexSet}; hence $g_i$ is incident with $T$ in the
sense of \noderef{TreeIncidentEdge}.  Equation (1) preserves the exact
color of $g_i$.  The same augmentation supplies a color missing at some
vertex of $A_i$ after expanding \noderef{TreeMissingColors}; that vertex
also lies in $V(T)$, so (2) preserves this missing-color witness.  Its
length and initial-segment clauses do not mention a coloring.  Thus all
clauses of
\noderef{TashkinovAugmentation} hold for the same ordered step under
$\sigma_n$.  Induction over the finite sequence, using
\noderef{TashkinovAugmentationSequence}, proves
\[
 \begin{cases}
  T_n\longrightarrow T,&\text{under $\sigma_n$ in the PE case},\\
  T_n+f_n=A\longrightarrow T,&\text{under $\sigma_n$ in the SE/RE cases}.
 \end{cases}                                                   \tag{6}
\]

The initial-segment and strict-extension premises for $T_n\subset T$
come directly from \noderef{ExtendedTashkinovTree}.  In PE, the first
route in (6) is exactly the PE eligibility alternative of
\noderef{ModColoring}.  In SE or RE, (4), $F_n=\{f_n\}$, and the second
route in (6) are exactly its other alternative.  Applying the universal
clause of (3) now supplies an actual
\noderef{CorrespondingExtendedTashkinovTree} certificate on the same
terminal tree $T$.  By \noderef{MaximumPropertyTransport}, the given
maximum property transfers to its replacement ETT.  This proves (ii).
Retain the PE augmentation sequence in (6), since it is an explicit
hypothesis needed below.

We construct the transported hierarchy, separating the PE and non-PE
cases.  In the PE case, all hypotheses of
\noderef{HierarchySplitter} have now been named: the lower-rung package,
maximum property and positivity, properness and exact one-uncoloredness
of both colorings, the PE combined-stability conclusion of
\noderef{HierarchyOutsideKempeStable}, and the first $\sigma_n$-sequence
in (6).  That lemma supplies a corresponding ETT with maximum property,
a hierarchy base using the same $R_n$ and starred core, and a hierarchy
with the same stages and divider positions.

Suppose instead that the terminal step is SE or RE.  Use the corresponding
ETT already obtained from \noderef{ModColoring}.  The definition
\noderef{CorrespondingExtendedTashkinovTree} gives equality of the rung
numbers, equality of every positive ladder tree (in particular the two
copies of $T_n$), equality of the raw stored sets and connecting-edge
sets, and equality of the extension kinds.  It does not assert equality
of internal PE or RE support data, and no such equality is used here.

Rebuild \noderef{HierarchyBase} for the replacement ETT as follows.
Take its core to be its terminal ladder tree $T_n$, transport the old
positive support position across the equality of the two copies of
$T_n$, and copy the unused closure field.  The terminal missing-color
field is reflexive because both the base coloring and the replacement
ETT's generating coloring are $\sigma_n$.  The parallel implication has
a false premise: equality of the terminal extension kinds and the present
SE/RE assumption rule out PE.  The nonparallel field is the defining
choice that the core is $T_n$.  Thus every field of the new base has been
specified.  For the old base, its nonparallel clause gives the same
identity $\mathcal B.\mathit{core}=T_n$.

Now construct a \noderef{TashkinovHierarchy} by copying the old divider
count, every ordered stage, and every divider position.  Its base-stage
field follows from the old base-stage equality and the corresponding
certificate's equality of the two terminal ladder trees.  Its terminal
stage remains literally $T$.  Every nested initial-segment relation,
every strict length inequality, positivity and increasing order of the
divider positions, and every divider-prefix identity is the corresponding
old field because the ordered stages and terminal tree were copied.  Its
core-inside-first field is the old one after rewriting both cores as
$T_n$.  This verifies every field of the hierarchy definition.

The two non-PE presentations satisfy
\noderef{SameHierarchyStages} field by field: their cores are both
$T_n$; their divider counts, stage families, and divider positions are
the copied data; equality of terminal kinds makes both PE propositions
false; and the conditional closure-equality clause is therefore vacuous.
Together with the PE output of \noderef{HierarchySplitter}, this gives in
all cases a corresponding maximum ETT, a hierarchy base, and a hierarchy
with the same paper-relevant stages and dividers as the original.

Apply \noderef{OutsideKempeGoodHierarchy} to these two presentations.
Its rung-equality and raw-stored-set hypotheses are precisely fields of
the chosen \noderef{CorrespondingExtendedTashkinovTree}; its same-stage
hypothesis was just proved, and its remaining hypotheses are the given
good hierarchy, properness, one-uncoloredness, and outside recoloring.
It supplies a good hierarchy under $\sigma_n$ and proves the relation
\noderef{SameGoodHierarchy}.  Consequently the corresponding
$D_{n,j}$ sets are equal and the Gamma values agree for every color in
each such common set.  This is exactly assertion (iii), including the
paper's meaning of ``the same Gamma sets.''

It remains to prove the final implication.  Put
\[
 C=\bigcup_{\eta\in D_{n,q+1}}\Gamma^q_\eta,
 \qquad
 D_{n,q+1}=\noderef{HierarchyTerminalStoredColors}.
\]
First we identify the control set appearing after subtraction.  By
\noderef{TreeMissingColors}, a color is in
$\overline{\varphi_n}(T)$ exactly when it is missing at some
$v\in V(T)$, and similarly for $\sigma_n$.  Applying (2) to that same
vertex in each direction proves
\[
       \overline{\varphi_n}(T)=\overline{\sigma_n}(T),
 \quad\text{hence}\quad
       \overline{\varphi_n}(T)-C
        =\overline{\sigma_n}(T)-C.                            \tag{7}
\]

Next fix a color $c$.  The set \noderef{BoundaryEdges} depends only on
$G$ and $V(T)$, so it is common to the two colorings.  If $f$ is one of
these exact boundary-edge identities, that definition supplies an
endpoint $u\in V(T)$ incident with $f$; hence
\noderef{TreeIncidentEdge} says that $f$ is incident with $T$, and (1)
gives $\sigma_n(f)=\varphi_n(f)$.  Filtering the common boundary by the
condition that the exact edge value is $c$, as in
\noderef{BoundaryColorEdges}, therefore gives
\[
 \partial_{\varphi_n,c}(V(T))
   =\partial_{\sigma_n,c}(V(T))                              \tag{8}
\]
as sets of edge identities, for every $c$.

Assume that $T$ is $C^{-}$-closed under $\varphi_n$.
By \noderef{TreeColorSetMinusClosed} and
\noderef{TreeColorSetClosed}, this says that for every
$c\in\overline{\varphi_n}(T)-C$, the left side of (8) is empty.
Equation (7) makes the quantified color set exactly
$\overline{\sigma_n}(T)-C$, and (8) makes the corresponding right side
empty.  A second application of \noderef{TreeColorSetClosed} and
\noderef{TreeColorSetMinusClosed} proves that $T$ is $C^{-}$-closed
under $\sigma_n$.  This proves (iv) and completes the lemma.
\end{proof}
